% REVIEW DRAFT: not an input of the current published manuscript.
% This file is not proof certification; see geometry-audit-2026-11-21.json.
\subsection{Flat degree and saturation for every principal-signature fibre}
\label{subsec:sr-general-flat-degree-saturation}

Here a host retains its full definition: \(E\) is CM,
\(E/E^+\) is unramified at every finite place, \(d\) is a prime
with \(d\equiv7\pmod {16}\), and
\[
 \varrho^2=-d,\qquad w_0=(1+\varrho)/2\in\mathcal O_E,\qquad
 \bar w_0=1-w_0.
 \tag{FD1}
\]
Put \(K_d=\mathbb Q(\varrho)\), write \(\mathfrak D_E\) for the
absolute different and \(R_E=|\operatorname{disc}E|^{1/[E:\mathbb Q]}\),
and set
\[
 x_d=\gamma+\log(2\pi/\sqrt d)-L'(1,\chi_d)/L(1,\chi_d),\qquad
 \Omega_E=x_d-\log R_E+\log\sqrt d .
 \tag{FD2}
\]
Thus \(\Omega_E=\operatorname{Re}(\mathrm{CM.logOm}(E,\chi_d))\).
The rank-one input is \eqref{eq:sr-cm-period-exact}, with metric
\eqref{eq:sr-cm-tangent-norm}.  It will only be applied to CM
elliptic curves; the passage to general fibres is proved below.
No abelian-Galois hypothesis on \(E/\mathbb Q\) is used.

Every degree is normalized by the number-field degree.
For logarithmic lengths over \(k\) the convention is
\[
 \deg N=-\frac1{[k:\mathbb Q]}
 \left(\sum_{\sigma:k\hookrightarrow\mathbb C}N_\infty(\sigma)
       +\sum_p\sum_{j:k\hookrightarrow\mathbb C_p}N_p(j)\right).
 \tag{FD3}
\]
We use \(|p|_p=p^{-1}\) and count all complex embeddings,
including both members of a conjugate pair.

Let \(\Lambda\) be a projective \(\mathcal O_E\)-lattice of rank
\(r=p+q\), with perfect alternating integral trace pairing and
Rosati compatibility.  Its period domain has signature \(p,q\)
at embeddings over \(\varrho\mapsto\sqrt{-d}\), and opposite
signature at the conjugates.  Fix a geometric component \(U\)
of a fine algebraic PEL moduli scheme, and write
\(V_\tau=\operatorname{Lie}(\mathcal A)_\tau\).
An algebraic frame function \(\phi\) is nonzero and satisfies
\[
                  \phi(Av,Bw)\det A=\phi(v,w)\det B
 \tag{FD4}
\]
on frames of \(V_\tau,V_{\bar\tau}\).
It is a generator of
\(\det V_\tau\otimes(\det V_{\bar\tau})^{-1}\).
For an integral finite-index \(E\)-basis \(e\) of the period
lattice, flatness means
\[
                   \phi(v,w)\det\binom C F=\rho\in\mathbb C^\times
 \tag{FD5}
\]
for the two period coefficient matrices.  A global free
\(\mathcal O_E\)-basis is not assumed.

For \(j:E\hookrightarrow\mathbb C_p\), put
\[
 D_j=\log\sup_{b\in\mathfrak D_E}|j(b)|_p,\qquad
                         h_{p,j}=\log|d|_p .
 \tag{FD6}
\]
Suppose its Hodge lengths and its lengths on integral eigenframes
of every semiabelian trait realization are constant on \(U\):
\[
 N_\infty(\sigma)=\log|\phi(v,w)|_\sigma
       +\tfrac12\log|\det\Gamma_v|_\sigma
       -\tfrac12\log|\det\Gamma_w|_\sigma,\qquad
 \widetilde N_p(j)=\log|\phi(v,w)|_j .
 \tag{FD7}
\]
Here the second formula is evaluated on integral eigenframes.
The source's corrected \(\mathrm{HasNorms}\) system is
\[
                    N_p(j)=\widetilde N_p(j)
                                        -q(D_{j|E}-h_{p,j}).
 \tag{FD8}
\]
This codifferent correction is essential.

\begin{srproposition}[Full flat degree]
\label{prop:sr-general-flat-degree}
For every host, signature and algebraic flat frame just described,
if the constant lengths exist, then they have finite support and
\[
 \deg N=\frac{p-q}{2}\Omega_E+q(\log d-\log R_E),\qquad
                  \deg\widetilde N=\frac{p-q}{2}\Omega_E .
 \tag{FD9}
\]
No field-degree, level, rank or isogeny-degree bound enters.
\end{srproposition}
\begin{proof}
We justify the CM evaluations on an arbitrary component.
At a rational prime \(p_0\), put
\(A=\mathcal O_E\otimes\mathbb Z_{(p_0)}\).
The different has a real generator \(\delta\) over \(A\):
the relative different of \(E/E^+\) is the unit ideal, and
ideals in the semilocal Dedekind ring of \(E^+\) are principal.
A denominator prime to \(p_0\) puts this generator in the
original ideal.
Multiplying the skew-hermitian trace pairing by \(\delta\)
makes it unimodular over \(A\).

It admits an orthogonal local basis with unit diagonal, including
above \(2\).  Indeed at each maximal ideal some \(a\) has
\(a-\bar a\) nonzero in its residue map: conjugation either
exchanges the two split factors or acts nontrivially on the
unramified quadratic residue extension.
If all diagonal values vanished, the values on \(u+av\)
would then force every cross pairing to vanish, contradicting
nondegeneracy.  Choose a vector of unit norm at every maximal
ideal by the Chinese remainder theorem, split off its orthogonal
line, and induct.

Perturb the first \(p\) vectors, preserving their congruence
modulo an arbitrarily large power of \(p_0\), so that their
whole \(p\)-by-\(p\) Gram block is positive definite at every
archimedean embedding over the chosen square root.
Choose a target positive-definite \(p\)-frame at each embedding;
scale inside the open homogeneous cone of such Gram blocks
and round in the prescribed integral lattice coset.
This gives the integral perturbation with the required
nonarchimedean congruences.
Their Gram determinant retains its local unit norm.
Orthogonal projection is integral at \(p_0\) since this block
is unimodular.  Its complement is negative by the signature.
Clearing denominators prime to \(p_0\) gives an integral
\(A\)-basis \(x=(x_+,x_-)\) with pure orthogonal blocks.
The associated pure period point remains in the same domain.
Its complex abelian variety is isogenous to
\[
             C_+^{\,p[E:K_d]}\times C_-^{\,q[E:K_d]} .
 \tag{FD10}
\]
Here \(C_\pm\) have the two CM types of \(K_d\).
For algebraicity, form the Serre tensor products with the
integral \(K_d\)-lattices in the two rational summands and quotient
by the finite subgroup specifying \(\Lambda\).
The polarization is the original integral trace pairing.
The CM curves, the subgroup, the action and the polarization
descend to a number field.  After finite extension the CM
curves have good reduction and the isogenies extend to their
semiabelian identity models.  Thus this constructs an algebraic
CM point of the actual component.

Put \(G=\det(\langle e_i,e_j\rangle_E)\ne0\), and write
\(x=T_{p_0}e,\ t_{p_0}=\det T_{p_0}\).
Self-duality of the alternating trace lattice gives
\[
 \log|jt_{p_0}|+\log|j\bar t_{p_0}|+\log|jG|
                                    =-(p+q)D_j .
 \tag{FD11}
\]
In fact its hermitian dual is
\(\mathfrak D_E^{-1}\Lambda\); take determinants in a local
basis.  At primes where \(e\) is a local basis the transition
determinants have norm one; only finitely many primes fail.

Let \(z_\pm\) be the periods of
Section~\ref{sec:sr-quadratic-period-witness}.
Evaluating (FD5) at one pure algebraic CM point gives
\[
 r_0z_+^p=\rho z_-^q,\qquad
 r_0=\phi(y,v,w)\,
       \tau(\det T_{p_0})/\tau(\det G_{x_-})\ne0 .
 \tag{FD12}
\]
All factors defining \(r_0\) are algebraic.
Enlarge the coefficient field \(k\) to contain them and the
CM curves.  The same \(r_0\) works at all the pure local points,
by the first identity and the fixed \(\rho\).

For the squared Hodge norms \(n_\pm(\sigma)\) of the integral
tangent generators of \(C_\pm\), the elliptic CM formula gives
\[
 -\frac1{[k:\mathbb Q]}\sum_\sigma\log n_+(\sigma)
 =-\frac1{[k:\mathbb Q]}\sum_\sigma\log n_-(\sigma)=x_d .
 \tag{FD13}
\]
Taking Gram determinants in (FD12) gives at every embedding
\[
 \begin{split}
 N_\infty(\sigma)={}&\log|\sigma r_0|
 +\frac p2(\log n_+(\sigma)+\tfrac12\log d)\\
 &-\frac q2(\log n_-(\sigma)+\tfrac12\log d)
                              +\tfrac12\log|\sigma(\tau G)|.
 \end{split}
 \tag{FD14}
\]
Indeed the pure Gram blocks are
\(n_\pm(\pm\tau(\varrho))\tau(G_{x_\pm})\), and
\(|\tau(\varrho)|=\sqrt d\).
Conjugating the algebraic CM point proves this at every embedding.

The finite calculation, with \(j:k\hookrightarrow\mathbb C_{p_0}\),
is
\[
 \widetilde N_{p_0}(j)=\log|jr_0|-\log|j(\tau t_{p_0})|-qD_{j|E}
                 +(q-p)(D_{j|E}-\log|j(\tau\varrho)|).
 \tag{FD15}
\]
To check its integrality factors, take a \(\mathbb Z\)-basis
\(b\) of \(\mathcal O_E\) and its trace-dual basis
\(b^*\subset\mathfrak D_E^{-1}\).
The \(\tau\)-coefficient in a Serre tensor product is extracted
by \(\sum_b\tau(b^*)b\).  Its integral lattice therefore uses
the codifferent.
Since the different of \(K_d\) is \((\varrho)\), the scalars
\(\delta_\pm\) changing pure frames to integral eigenframes satisfy
\[
       \log|j\delta_+|=\log|j\delta_-|
                           =D_{j|E}-\log|j(\tau\varrho)|.
 \tag{FD16}
\]
Apply the trace-dual coefficient map to a local basis; all
remaining matrices and their inverses are integral, hence have
unit determinant.  The negative Gram block contributes
\(-qD_{j|E}\).
Now (FD4), (FD12) and (FD16) prove (FD15).
This uses a rational projector to compute the lattice; it does
not suppose that spectral projectors are integral.

As \(2\log|j(\tau\varrho)|=\log|d|_{p_0}\), (FD8) simplifies to
\[
 N_{p_0}(j)=\log|jr_0|-\log|j(\tau t_{p_0})|
                     -(p+q)D_{j|E}+\frac{p+q}{2}\log|d|_{p_0}.
 \tag{FD17}
\]
This proves finite support, supported only on \(r_0,d,\mathfrak D_E\)
and the finite-index frame \(e\).
Conjugation permutes the \(p_0\)-adic embeddings, so summing
(FD11) replaces the sum of \(-\log|jt_{p_0}|\) by half the sum
of \(\log|jG|+(p+q)D_j\).
The normalized finite sums of \(D_j\) and \(\log|d|_p\) are
\(-\log R_E\) and \(-\log d\).
The product formula cancels \(r_0\) and \(G\).  Consequently
(FD3), (FD13)--(FD14) give
\[
       \deg N=\frac{p-q}{2}x_d-\frac{p+q}{2}\log R_E
                                      +\frac{p+3q}{4}\log d .
\]
This is the first formula in (FD9); the correction (FD8)
changes the degree by \(q(\log R_E-\log d)\) and proves the second.
Normalized embedding sums and degrees are invariant under
finite extension, so the result descends from \(k\).
\end{proof}

\paragraph{Actual integral models and polarized uniformization.}
The precise classical background used for frame existence is
the complete-normal-ring equivalence of polarized
degenerations with Mumford data, including its cubical-line
and endomorphism variants: Lan, \S4.1, Theorems~4.4.16 and~5.1.1.4.
The base may be any noetherian normal domain complete for a
radical ideal.  Proper integral PEL models, including absolute
ramified primes, are supplied by the normalized toroidal models
and extended semiabelian families in Lan's ramified-characteristic
Theorems~11.2 and~11.4.
See the author's
\href{https://www.kwlan.org/articles/cpt-PEL-type-thesis-revision.pdf}
{Arithmetic compactifications} and
\href{https://www.kwlan.org/articles/cpt-ram-flat.pdf}
{Compactifications in ramified characteristics}.
We verify the extra outputs required here.

Choose level divisible by two distinct odd primes split in \(K_d\),
with complementary ideals \(\mathfrak M,\bar{\mathfrak M}\).
Increase the powers of both coprime auxiliary levels so that each defines a neat principal-congruence subgroup.  Then the prime-to-\(p\) level image is neat for every \(p\), and its factor is at least \(3\).  Add full auxiliary level as well.
After finite extension the given point has that level.
Take the normal closure of its geometric PEL component in a
projective toroidal compactification.
It is proper and flat over the coefficient valuation ring
and carries the extended semiabelian family and action.
The actual symmetric cubical line is fixed by the polarization
and \(w_0\), by pulling the biextension back under \((1,w_0)\);
the primitive identity in (FD1) identifies it with the original
rigidified line.  Thus the model carries the same line, not
just a polarization isogeny.
After separating the constant field its special fibre is
connected by Stein factorization.

Let \(R\) be a complete noetherian normal local characteristic-zero
domain mapping to this model, with nonempty abelian generic locus
and split torus in its closed fibre.
The degeneration equivalence gives an abelian \(R\)-scheme \(B\),
a Raynaud extension \(1\to T_X\to\widetilde G\to B\to1\), a period
lattice \(Y\), a principal polarization \(Y\simeq X\), the
extension maps, biextension trivialization \(\tau\), and cubical
trivialization \(\psi\).
The formal and generic quotients identify with the given family.
Functoriality transfers the action and Rosati involution.
The rigidified primitive identity gives
\[
 (1+w_0)^*M\otimes M^{-2}\otimes w_0^*M^{-1}\otimes e^*M^2
     \simeq\mathcal O_B,\qquad
                 \psi(y)=\tau(y,w_0y)\,u(c(y)),
 \tag{FD18}
\]
with rigidifications understood.
This is the actual trace-form-twist identity; its scalar is
fixed by rigidification.

On every toroidal chart the period trivializations are units
times boundary monomials.  Each boundary monomial is invertible
on the abelian locus.  Therefore every \(\psi(y)\), in each
line trivialization, is a regular unit on every affine open
of that locus.  This proves the source's \(\mathrm{PsiUnitOn}\)
clause, not only nonvanishing at the generic point.

For every \(m\ge1\) and algebraically closed
\(F\supset\operatorname{Frac}R\), the torsion comparison is
\[
 \frac{\{(y,g):y\in Y,\ g\in\widetilde G(F),\
                                  [m]g=\iota(y)\}}
      {(y,g)\sim(y+mz,g+\iota(z))}
                                      \simeq A[m](F).
 \tag{FD19}
\]
The quotient construction defines the map.
Its kernel is exactly the displayed period relation, and lifting
an \(m\)-torsion point proves surjectivity.
One can first check this over a completed algebraic closure,
where it is the analytic quotient, and then descend the finitely
many roots of a period and torsion points through the algebraic
cubical construction.  These formulas commute with field
extension.  Addition and the action are preserved.
The Weil pairing on both sides is the commutator of the two
lifts to the cubical theta extension, namely the ratio of the
two biextension trivializations.  Hence it too is preserved,
also when \(m\) is divisible by the residue characteristic.
The construction verifies all of \(\mathrm{Model.N3}\),
including (FD18), the unit clause and (FD19).
The independent arc-uniqueness condition \(\mathrm{N4}\) is not
used in this subsection.


\paragraph{The full integral Hodge lattice.}
Formal coordinates and the rigidified cubical line construct
a rank twice-dimension bundle \(\mathcal H\) with
\[
 0\to\operatorname{Lie}(\mathcal G)^\vee
         \to\mathcal H\to\operatorname{Lie}(\mathcal G)\to0 .
 \tag{FD20}
\]
Here and below its left action is the conjugate dual action.
A parameter change and change of cubical trivialization act by
\[
                  (v,f)\longmapsto(Cv,C^{-t}(f+Z^tv)).
 \tag{FD21}
\]
The cocycle law proves gluing.
For \(a\in\mathcal O_E\), let \(R_a\) be its tangent matrix
and \(A_a\) the mixed derivative of its cubical comparison.
The action on this bundle is
\[
                 (v,f)\longmapsto
                     (R_av,R_{\bar a}^{\,t}f-A_a^{\,t}v).
 \tag{FD22}
\]
Differentiating the cubical and Rosati identities proves
additivity and multiplicativity.
Over the characteristic-zero abelian locus it is the full
de Rham realization, in its covariant convention.
Its characteristic polynomial at \(a\) is
\(\operatorname{char}(a\mid E)^r\), since the conjugate
Lie multiplicities add to \(r=p+q\).
This polynomial identity extends over the reduced normal model.

We need actual local freeness over
\(\mathcal O_E\otimes\mathcal O_{\rm model}\).
At a geometric special point the reduction of polarized
Mumford data gives a torus extension of an abelian variety
\(B\), and the actual line is pulled back from a principally
polarized line on \(B\).
The torus character lattice is a projective
\(\mathcal O_E\)-lattice.

The full Hodge module of \(B\) over a perfect field \(k\)
of characteristic \(p\) is the reduction of its integral
Dieudonne realization, with the conjugate-dual convention
fixed by its polarization.
We use only the standard functorial identification with
degree-one crystalline/de Rham realization, not an isogeny
comparison with denominators; see
\href{https://numdam.org/book-part/AST_1979__63__17_0/}
{Berthelot--Messing, \emph{Theorie de Dieudonne cristalline I}},
and the abelian comparison in Berthelot--Breen--Messing,
\emph{Theorie de Dieudonne cristalline II}.
The integral module is free over \(W(k)\).
Each factor of \(\mathcal O_E\otimes W(k)\), after separating
unramified embeddings, is a discrete valuation ring.
A power of its uniformizer is \(p\) times a unit; hence the
module is torsion free over this ring, and so free.
The total characteristic polynomial fixes its ranks.
This reasoning includes \(p=2\) and ramified local number fields.

There is an elementary check of the torus pieces.
Let \(P:\operatorname{Lie}G\to\operatorname{Lie}B\) be the
quotient and \(T=\ker P\).
A cubical change of gauge makes
\(A_a=P^tA_{B,a}P\) in (FD22).
The map \((v,f)\mapsto f|_T\) is equivariant.
On its kernel, a linear section \(Q\) of \(P\) gives
\[
                         (v,f)\longmapsto(Pv,Q^tf).
\]
It is an equivariant surjection onto the full Hodge module
of \(B\), with kernel \(T\).
For instance \(f|_T=0\) says \(f=P^tQ^tf\), which verifies
both the mixed term and the adjoint action in this formula.
Thus the full Hodge module has three successive pieces:
\(T\), the full Hodge module of \(B\), and the conjugate
dual of \(T\).
They are free over the residue algebra.
Extensions with free quotient split as modules over that
algebra, so the full module is free.
The characteristic polynomial identity gives rank \(r\)
on each factor.

Lift such a residue basis to a local ring of the model.
It gives a map
\((\mathcal O_E\otimes\mathcal O_{\rm model})^r\to\mathcal H\).
Nakayama gives surjectivity, and equal ordinary ranks over
the local base give injectivity.
Descent from a geometric residue field is faithfully flat.
This proves the required local freeness; characteristic
polynomials alone would not suffice over the nonreduced
residue algebra.

\begin{srlemma}[All character quotient degrees]
\label{lem:sr-flat-fibre-character-degree-zero}
On every proper curve in a geometric special fibre, every
residue character \(\chi:\mathcal O_E\to k\) has a locally
free rank-\(r\) quotient
\[
 \mathcal Q_\chi=\mathcal H/
          \sum_{a\in\mathcal O_E}(a-\chi(a))\mathcal H
                       \quad\text{with}\quad\deg\mathcal Q_\chi=0 .
 \tag{FD23}
\]
\end{srlemma}
\begin{proof}
Normalize the curve and extend its constant field when necessary;
degree zero descends.
The formal Hodge construction in characteristic \(p\) gives
the Hodge and conjugate exact sequences
\[
 \begin{split}
 0&\to\mathcal L^{\vee,\bar{\ }}\to\mathcal H\to\mathcal L\to0,\\
 0&\to F^*\mathcal L\to\mathcal H
                         \to F^*(\mathcal L^{\vee,\bar{\ }})\to0,
 \end{split}
 \tag{FD24}
\]
where \(\mathcal L=\operatorname{Lie}\mathcal G\).
These are assertions about the full integral lattices.
For abelian fibres they are the Hodge and Cartier sequences.
Their formal-coordinate proof also works for the torus
extension with a pulled-back principal line:
if \(H,A,S\) are the Hasse--Witt and mixed cubical matrices,
the second sequence has maps
\[
              u\mapsto(Hu,A^tu),\qquad
              (v,f)\mapsto(A-H^tS)v-H^tf.
\]
The cubical identities give zero composite.
The first matrix has full column rank modulo every closed
valuation ideal: the torus Hasse--Witt block is invertible,
and on the abelian quotient its possible common kernel is
the infinitesimal translation kernel of the principal line,
which is zero.
A full-rank minor is a unit; the displayed adjoint identity
then gives split exactness over the local valuation ring.
The coordinate changes (FD21) transport these exact sequences.
The actions on their Frobenius terms are Frobenius twists,
as is checked by applying \(a\) before or after Frobenius.

The finite image of \(\chi\) is contained in a finite field
\(k_0\subset k\).  Put \(A_0=k_0\otimes_{\mathbb Z}\mathcal O_E\).
For an idempotent \(e\in A_0\), set \(d(e)=\deg(e\mathcal H)\).
The two exact sequences and
\(\deg F^*\mathcal V=p\deg\mathcal V\) imply
\[
                              d(F(e))=p\,d(e),
 \tag{FD25}
\]
where \(F\) raises only the \(k_0\)-coefficients to their
\(p\)-th powers.  A positive iterate of \(F\) is the identity,
so \(d(e)=p^m d(e)\) and \(d(e)=0\).

Let \(e_\chi\) be the idempotent of the local Artin factor
of \(A_0\) belonging to \(\chi\).
Filter that factor by ideals with one-dimensional residue-field
successive quotients.
Since \(\mathcal H\) is locally free over \(A_0\), this
filtration remains exact on the local integral lattices.
Multiplication by a representative of a simple ideal
quotient identifies the corresponding quotient bundle with
\(\mathcal Q_\chi\).
Its possible change is a constant unit.
If \(m_\chi>0\) is this Artin factor's length, additivity gives
\[
                       0=d(e_\chi)=m_\chi\deg\mathcal Q_\chi.
\]
This proves degree zero.
Applying the character to a free rank-\(r\) module proves
local freeness.  In particular no integral projector onto
a colliding characteristic-zero embedding was used.
\end{proof}

\paragraph{Ramified determinants and the actual finite norm.}
We still must compare this quotient with literal integral
eigenspaces of the Lie lattice.
Work in an algebraically closed complete valued field \(C\)
of characteristic zero, with valuation ring \(V\).
Let \(S\) be the residue cluster of an embedding \(\tau\).
Choose \(\pi\in\mathcal O_E\) generating its ramified local
factor over the maximal unramified subring.
Approximation to the finitely many local uniformizers gives
such a global element.
After separating residue embeddings, the evaluation algebra is
\(V[T]/\prod_{s\in S}(T-a_s)\), \(a_s=s(\pi)\).
Write \(r_s=\operatorname{rk}\mathcal L_{\bar s}\) and \(r=p+q\).
Its fixed determinant factor is
\[
 \gamma_\tau=
  \frac{\prod_{s\ne\tau}(a_\tau-a_s)^{r_s}}
       {\left(\prod_{s\ne\tau}(a_\tau-a_s)\right)^r}\ne0 .
 \tag{FD26}
\]
It depends on embeddings and their ranks, not on the point.

Here is the local calculation.
Take a free \(\mathcal O_E\otimes V\)-frame \(t_1,\ldots,t_r\)
of the full Hodge lattice.
In its \(\tau\)-component choose an adapted basis \(u\):
the first \(p\) vectors project to a saturated integral
basis \(v\) of \(\mathcal L_\tau\), the others are dual to
an integral basis \(w\) of \(\mathcal L_{\bar\tau}\).
If the \(\tau\)-projection of the free frame is \(A u\), then
\[
                              \det A\in V^\times\gamma_\tau .
 \tag{FD27}
\]
We prove this using the following exact lattice computation.

Write \(R=V[T]/\prod_{s\in S}(T-a_s)\) and
\(\mathscr M=R^r\subset\bigoplus_s C^r\), using evaluation.
Let \(g_s:C^r\to W_s\) be the projection of the full Hodge
space to its \(s\)-Lie component, and put
\[
 F_s=\ker g_s,\quad
 \mathscr F=\mathscr M\cap\bigoplus_sF_s,\quad
 \mathscr L=\mathscr M/\mathscr F,\quad
 \Delta_\tau=\prod_{s\ne\tau}(a_\tau-a_s).
 \tag{FD27a}
\]
The images \(W_s\) are independent, so \(\mathscr L\) is
exactly the integral Lie image of this residue cluster.
The module \(\mathscr F\) is saturated and \(T\)-stable.
Polynomial evaluation gives
\[
 e_\tau\mathscr M=V^r,\qquad
 \ker(T-a_\tau\mid\mathscr M)=\Delta_\tau V^r,\qquad
 \operatorname{coker}(T-a_\tau\mid\mathscr M)=V^r .
 \tag{FD27b}
\]
For the middle identity, a polynomial supported only at
\(a_\tau\) is a multiple of
\(\prod_{s\ne\tau}(T-a_s)\); division by these monic factors
takes place over \(V\).

Here are all index calculations, valid also for a nondiscrete
valuation.  For a finitely presented torsion \(V\)-module
\(Z\), let \(\lambda(Z)\) be the sum of the valuations of its
nonzero Smith diagonal entries.
Every finite matrix over a valuation ring admits such a
diagonalization by pivoting at an entry of least valuation.
Thus this length is well defined and additive in exact
sequences of these modules.  In the discrete case it is the
usual length times the valuation of a uniformizer.
All modules here are kernels, images or cokernels of finite
matrices; the saturated kernels and projected images are
finite free.  This follows from the same finite elimination,
so no noetherian hypothesis on \(V\) is being used.

Set \(D=T-a_\tau\), and define
\[
 K=\ker(D\mid\mathscr F),\quad P=e_\tau\mathscr F,\quad
 W=\mathscr F\cap\bigoplus_{s\ne\tau}F_s,\quad
 Z=\operatorname{coker}(D\mid\mathscr F)_{\rm tors}.
\]
Projection gives the exact sequence
\[
                      0\to W\to\mathscr F/K\to P/K\to0 .
 \tag{FD27c}
\]
Indeed its kernel is exactly the elements with zero
\(\tau\)-component.  The quotient \(P/K\) is torsion.
Moreover \(D(\mathscr F)\subset W\), and
\(\mathscr F/W=P\) is torsion free, hence
\(Z=W/D(\mathscr F)\).
The map \(D\) induces an injection \(\mathscr F/K\to W\).
It follows from (FD27c) that
\(\lambda(D(\mathscr F)/D(W))=\lambda(P/K)\).
The determinant of \(D\) on \(W\) is
\(\prod_{s\ne\tau}(a_s-a_\tau)^{r_s}\), because its
rational eigenspaces are exactly the \(F_s\).
Consequently
\[
        \lambda(Z)+\lambda(P/K)
                  =\sum_{s\ne\tau}r_s\,v(a_\tau-a_s).
 \tag{FD27d}
\]

Apply the kernel--cokernel sequence for \(D\) to
\(0\to\mathscr F\to\mathscr M\to\mathscr L\to0\).
By the torsion-free cokernel in (FD27b), it gives
\[
 0\to K\to\Delta_\tau V^r
          \to\mathscr L\cap W_\tau\to Z\to0 .
 \tag{FD27e}
\]
To check the last surjection explicitly, an element of \(Z\)
has zero rational \(\tau\)-component; its image in
\(\operatorname{coker}(D\mid\mathscr M)\) is torsion and
therefore zero, so it lifts to \(D m\).
The class of \(m\) lies in the indicated integral
\(\tau\)-Lie module.
Conversely every such class maps to a torsion cokernel class.
Thus the last term is precisely \(Z\), not an unspecified
defect bounded by it.

Choose a basis \(u_L\) of \(\mathscr L\cap W_\tau\), with
arbitrary lifts to \(C^r\), and a basis \(u_P\) of \(P\).
Use \(K\) first in (FD27e), and then enlarge it to \(P\).
The two determinant-index equalities give
\[
 \begin{split}
 v\det(u_L,u_P)
  &=r\,v(\Delta_\tau)-\lambda(Z)-\lambda(P/K)\\
  &=\sum_{s\ne\tau}(r-r_s)v(a_\tau-a_s).
 \end{split}
 \tag{FD27f}
\]
The second equality is (FD27d).  This is the complete
determinant identity; flag-dependent indices cancel.

For application to the full Hodge lattice,
\(\mathscr L\cap W_\tau\) has basis \(v\).
The projected kernel \(P\) identifies with the dual of
the saturated integral \(\bar\tau\)-Lie lattice:
projection of \(\mathcal L^\vee\) to this dual component
is restriction of linear functionals to
\(\mathcal L_{\bar\tau}\).
That submodule is saturated, hence a direct summand over
the valuation ring, so every integral functional extends.
Therefore the dual basis of \(w\) is a basis of \(P\).
The adapted frame \(u\) is precisely the appended frame
in (FD27f).  The projected free frame \(t\) is a basis of
\(e_\tau\mathscr M=V^r\), and \(t=A u\).
Taking inverse determinants in (FD27f) proves (FD27),
with exactly the factor (FD26).
Equality of valuations of two nonzero scalars means that
their ratio is a unit of \(V\), so the asserted unit
identity holds, not only an inequality.

If the Lie frames change by \(P,Q\), the comparison reads
\[
                   \det A\,\det Q
                             =u\gamma_\tau\det P,\qquad u\in V^\times .
 \tag{FD28}
\]
This equation survives specialization to an actual integral
trait point.  Its eigenlattices are the kernels of
\(av-\tau(a)v=0\), for a finite integral basis of
\(\mathcal O_E\); flat scalar extension preserves these kernels.
The mixed matrices (FD22) specialize by the same changes
of coordinates.  This explains why the quotient and the
actual saturated eigenlattices have exactly the constant
(FD26), even at absolute ramified primes.

On a normal proper model over a coefficient DVR \(O\),
put \(\mathscr Q=\det\mathcal Q_\tau\).
The rational section \(s\) coming from a flat generic frame
is a generator on the generic compactified component.
The preceding degree-zero lemma implies that there are
integers \(M>0,b\) with
\[
                          s^M/\varpi^b
               \text{ a generator of }\mathscr Q^{\otimes M}
                                  \text{ at every point}.
 \tag{FD29}
\]
Indeed its divisor is vertical.
Subtract the least ratio of a component coefficient to the
corresponding special-fibre multiplicity, and clear denominators.
The remaining divisor is effective and has coefficient zero
on some special-fibre component.
If nonzero, connectedness gives a zero-coefficient component
meeting its support.  A proper curve in that component,
through a meeting point and not contained in the support,
has positive degree against the effective divisor, contrary
to numerical triviality.
One obtains the curve by hyperplane sections on a projective
model and normalization.
The effective divisor is therefore zero; normality identifies
Cartier divisors from their codimension-one coefficients.
This proves (FD29).

Together (FD28)--(FD29) give, on actual integral trait eigenframes,
\[
              \log|\phi(v,w)|_p
                       =\frac bM\log|\varpi|_p+\log|\gamma_\tau|_p .
 \tag{FD30}
\]
Every term on the right is fixed on the component.
This is actual finite norm constancy.
The general CM evaluation (FD15) computes its value;
no estimate for the valuation of (FD26) is needed.

\begin{srproposition}[The actual adelic flat frame]
\label{prop:sr-actual-adelic-flat-frame}
For every nonzero principal-signature host datum, after
fine level and finite coefficient extension each geometric
period component has an algebraic flat frame with all
constants required in Proposition~\ref{prop:sr-general-flat-degree}.
The finite constants apply to every integral semiabelian
trait realization with the same generic point and cubical line.
\end{srproposition}
\begin{proof}
On the simply connected period domain use (FD5).
Its multiplier is the determinant of integral unitary
monodromy.  This is a unit whose absolute value at every
embedding is one, hence a root of unity.
Full level at least \(3\) makes it \(1\):
if \(\zeta\ne1\) and \(\zeta-1\in m\mathcal O_E\),
its norm is at least \(m^{[E:\mathbb Q]}\) and at most
\(2^{[E:\mathbb Q]}\).
Thus the determinant descends.
The polarized Gram identity normalizes its Hodge length to one.

This unit extends to the determinant line on a proper
characteristic-zero toroidal resolution.
In a Mumford boundary chart the monodromy logarithm occurs
in an off-diagonal block of the logarithmic de Rham transition.
This matrix is triangular with diagonal one.
Consequently its determinant has no boundary zero, pole or
logarithmic term.
The canonical extension of the determinant therefore agrees
with the determinant of (FD20).
The unit extends across the boundary divisors and, by normality,
the remaining codimension at least two sets.
Resolution is used only in characteristic zero.

GAGA makes the unit an algebraic line generator.
It descends to a number field: the line is defined there,
its section space commutes with extension to \(\mathbb C\),
and evaluation becomes an isomorphism after this faithfully
flat extension.
A proper birational map to a normal scheme pushes the
structure sheaf forward to the structure sheaf, so the
unit descends from the resolution to the original generic
component.  It gives (FD4)--(FD5).
Two such algebraic flat frames differ by a nonzero
coefficient-field scalar, since their ratio is a function
on a proper connected integral resolution.

At every complex embedding the same argument gives (FD7);
complex conjugation handles an embedding reversing the
chosen square root.
At a finite place use the actual normal proper model,
its full polarized uniformization (FD18)--(FD19), and
the full Hodge lattice.
Equations (FD23)--(FD30) give the constant on every trait.
Properness extends the point to this model.
Uniqueness of a semiabelian extension over a normal trait
identifies the resulting family and its cubical line
with the given realization, so this compares its actual
integral eigenframes.
Finite extensions separating the level or torus characters
do not change the conclusion.

Choose one algebraic flat frame.
At each other place its scalar ratio with the constructed
frame preserves constancy.
This gives a single system \(\widetilde N\) at all places.
Proposition~\ref{prop:sr-general-flat-degree} proves finite
support for it.  No uninstantiated global geometric package
remains in this construction.
\end{proof}

\begin{srtheorem}[General principal-signature flat saturation]
\label{thm:sr-general-flat-saturation}
Let \(L\) be any number field, and let
\(\mathcal G/\mathcal O_L\) be a semiabelian scheme
whose generic fibre is a principally polarized abelian
variety with its host action and the actual symmetric
rigidified cubical line specified by (FD1).
Give \(B=\operatorname{Lie}\mathcal G\) its actual Hodge metric.
Suppose \(\rho\in L,\ \rho^2=-d\), and \(p,q\ge0\) satisfy
after every algebraically closed characteristic-zero extension
\[
 \operatorname{rk}\operatorname{Lie}(\mathcal G)_s
       =\begin{cases}p&s(\varrho)=\rho,\\q&s(\varrho)=-\rho.
         \end{cases}
 \tag{FD31}
\]
Let \(s_1,s_2:\mathcal O_E\hookrightarrow\mathcal O_L\) be
conjugate under every complex embedding of \(L\).
Their integral eigenmodules \(B_{s_i}\) are saturated.
If \(s_1(\varrho)=\rho\), then
\[
               \widehat{\deg}_n B_{s_1}
                     -\widehat{\deg}_n B_{s_2}
                                   =\frac{p-q}{2}\Omega_E .
 \tag{FD32}
\]
In particular there is one absolute \(C\), chosen before
\(E,L,p,q\), such that
\[
 \left|\widehat{\deg}_n B_{s_1}-\widehat{\deg}_n B_{s_2}\right|
 \le C(\operatorname{rk}B_{s_1}+\operatorname{rk}B_{s_2})
       \left(1+\log R_E+|L'(1,\chi_d)/L(1,\chi_d)|\right).
 \tag{FD33}
\]
This is the full producer \(\mathrm{FlatSat}\).
\end{srtheorem}
\begin{proof}
If \(p+q=0\) both modules are zero.
Otherwise choose a complex embedding carrying \(\rho\)
to the fixed square root.
The actual Betti lattice and principal trace pairing of
the given fibre give precisely the datum considered above.
After finite extension and level it is a point of one
actual component.  Use its adelic flat frame.

Choose integral independent \(x_1,x_2\) spanning its two
generic eigenspaces, up to finite index in their saturated
lattices.  At infinity put
\(N_{i,\infty}=\tfrac12\log|\det\operatorname{Gram}(x_i)|\).
At a finite place choose an integral basis \(y_i\) of the
saturated local eigenlattice and write \(x_i=T_i y_i\);
a denominator of local norm one can first be cleared.
Put \(N_{i,p}=\log|\det T_i|_p\).
The determinant-ideal formula gives
\[
                         \deg N_i=\widehat{\deg}_n B_{s_i}.
\]
This includes nonprincipal determinant ideals, since the
indices of a finite-index integral frame are precisely the
local valuations of \(\det T_i\).

Let \(a=\phi(y,x_1,x_2)\ne0\), an algebraic scalar in the
coefficient extension.
The Gram identity at infinity and (FD4) at finite places give
\[
                       \widetilde N_v
                           =\log|a|_v+N_{1,v}-N_{2,v}
\]
at every place.  At finite places this uses the actual
semiabelian trait clause, not a good-reduction substitute.
Summing and applying the product formula to \(a\) yields
\[
              \deg\widetilde N
                       =\widehat{\deg}_n B_{s_1}
                                    -\widehat{\deg}_n B_{s_2}.
\]
Now (FD9) proves (FD32).
Finite flat scalar extension preserves the saturated
eigenmodules by their finite kernel description.
Hodge metrics pull back, and normalized determinant degrees
are invariant.  Hence this identity descends to \(L\).

Finally
\[
          \Omega_E=\gamma+\log(2\pi)-\log R_E
                                  -L'(1,\chi_d)/L(1,\chi_d).
\]
Use \(|p-q|/2\le p+q\) and \(R_E\ge1\).
This proves (FD33) with an absolute constant.
If \(s_1(\varrho)=-\rho\), swap the pair for the absolute-value
conclusion.
\end{proof}

\paragraph{Exact producer audit.}
The general CM computation is
\texttt{FlatDegree.flatDegree}, with the actual formulas
in \texttt{FrameFunctionCmValue},
\texttt{FrameFunctionCmHeight} and
\texttt{FrameFunctionCmNumberField}.
The adelic frame construction supplies the uses of
resolution, GAGA and integral models in \texttt{FlatZero},
\texttt{FlatFirstFromResolutionAndGaga} and
\texttt{FlatTheoremOfParts}.
The full Hodge lattice and Frobenius--Artin argument are
\texttt{HodgeLatticeDvrDegree} and
\texttt{SplitHermitianGeneratorsExist}.
The ramified determinant (FD26)--(FD30) is the actual
\texttt{CmTypeValuationEmbedding}/\texttt{EigenLinePAdicFrame}
comparison.
The final determinant-ideal comparison is
\texttt{ExistsCoordLattice} and \texttt{ExistsSmulMemSpan}.
Thus (FD32)--(FD33) give the substantive conclusions of
\texttt{flat\_hermitian\_fibre\_signature} and
\texttt{flat\_saturation\_of\_degree} in
\texttt{PhiDegFibLatLeOfIsOFamilyOfIsAdapted}.
The host's finite-unramified CM-extension hypothesis is retained.
There is no quadratic-CM-fibre restriction, abelian-Galois
assumption, rank-dependent constant or coefficient-field
degree loss.  The later integral-point commensurability
and global height budget are separate conclusions.
